\chapter{2004年上海卷（秋文）}
\section{填空题}
\begin{problem}
若 \(\tan \alpha = \frac{1}{2}\)，则 \(\tan \left(\alpha + \frac{\pi}{4}\right) = \fillinblank{}\).
\end{problem}

\begin{problem}
设抛物线的顶点坐标为 \( (2,0) \)，准线方程为 \(x = -1\)，则它的焦点坐标为 \fillinblank{}.
\end{problem}

\begin{problem}
设集合 \(A = \{5, \log_2(a + 3)\}\)，集合 \(B = \{a, b\}\). 若 \(A \cap B = \{2\}\)，则 \(A \cup B = \fillinblank{}\).
\end{problem}

\begin{problem}
设等比数列 \(\{a_n\}\) \((n \in \mathbf{N})\) 的公比 \(q = -\frac{1}{2}\)，且 \(\lim_{n\to\infty}(a_1 + a_3 + a_5 + \cdots + a_{2n-1}) = \frac{8}{3}\)，则 \(a_1 = \fillinblank{}\).
\end{problem}

\begin{problem}
设奇函数 \(f(x)\) 的定义域为 \([-5, 5]\). 若当 \(x \in [0, 5]\) 时，\(f(x)\) 的图象如右图，则不等式 \(f(x) < 0\) 的解是 \fillinblank{}.

\bitmapfigure[page=1,pagebox=mediabox,viewport=184 428 328 504,clip,max width=0.30\linewidth,max totalheight=4.0cm]{source/普通高考/2004/2004上海文.pdf}
\end{problem}

\begin{problem}
已知点 \(A(-1,5)\) 和向量 \(\bs{a} = (2,3)\)，若 \(\overrightarrow{AB} = 3\bs{a}\)，则点 \(B\) 的坐标为 \fillinblank{}.
\end{problem}

\begin{problem}
当 \(x, y\) 满足不等式 \(\begin{cases} 2 \leqslant x \leqslant 4, \\ y \geqslant 3, \\ x + y \leqslant 8 \end{cases}\) 时，目标函数 \(k = 3x - 2y\) 的最大值为 \fillinblank{}.
\end{problem}

\begin{problem}
圆心在直线 \(x = 2\) 上的圆 \(C\) 与 \(y\) 轴交于两点 \(A(0, -4), B(0, -2)\)，则圆 \(C\) 的方程为 \fillinblank{}.
\end{problem}

\begin{problem}
若在二项式 \((x + 1)^{10}\) 的展开式中任取一项，则该项的系数为奇数的概率是 \fillinblank{}.（结果用分数表示）
\end{problem}

\begin{problem}
若函数 \(f(x) = a|x - b| + 2\) 在 \([0, +\infty)\) 上为增函数，则实数 \(a, b\) 的取值范围是 \fillinblank{}.
\end{problem}

\begin{problem}
教材中“坐标平面上的直线”与“圆锥曲线”两章内容体现出解析几何的本质是 \fillinblank{}.
\end{problem}

\begin{problem}
若干个能唯一确定一个数列的量称为该数列的“基本量”. 设 \(\{a_n\}\) 是公比为 \(q\) 的无穷等比数列，下列 \(\{a_n\}\) 的四组量中，一定能成为该数列“基本量”的是第 \fillinblank{}组.（写出所有符合要求的组号）

\begin{circlelist}
\item \(S_{1}\) 与 \(S_{2}\);
\item \(a_{2}\) 与 \(S_{3}\);
\item \(a_{1}\) 与 \(a_{n}\);
\item \(q\) 与 \(a_{n}\).
\end{circlelist}

其中 \(n\) 为大于 \(1\) 的整数，\(S_{n}\) 为 \(\{a_{n}\}\) 的前 \(n\) 项和.
\end{problem}

\section{单选题}
\begin{problem}
在下列关于直线 \(l\)、\(m\) 与平面 \(\alpha\)、\(\beta\) 的命题中，真命题是（\quad）

\choices
    {若 \(l \subset \beta\) 且 \(\alpha \perp \beta\)，则 \(l \perp \alpha\)}
    {若 \(l \perp \beta\) 且 \(\alpha \parallel \beta\)，则 \(l \perp \alpha\)}
    {若 \(l \perp \beta\) 且 \(\alpha \perp \beta\)，则 \(l \parallel \alpha\)}
    {若 \(\alpha \cap \beta = m\) 且 \(l \parallel m\)，则 \(l \parallel \alpha\)}
\end{problem}

\begin{problem}
三角方程 \(2\sin \left(\frac{\pi}{2} - x\right) = 1\) 的解集为（\quad）

\choices
    {\(\left\{x\mid x = 2k\pi +\frac{\pi}{3},k\in \Z\right\}\)}
    {\(\left\{x\mid x = 2k\pi +\frac{5\pi}{3},k\in \Z\right\}\)}
    {\(\left\{x\mid x = 2k\pi \pm \frac{\pi}{3},k\in \Z\right\}\)}
    {\(\left\{x\mid x = k\pi +(-1)^k,k\in \Z\right\}\)}
\end{problem}

\begin{problem}
若函数 \(y = f(x)\) 的图象与函数 \(y = \lg (x + 1)\) 的图象关于直线 \(x - y = 0\) 对称，则 \(f(x) =\)（\quad）

\choices
    {\(10^{x} - 1\)}
    {\(1 - 10^{x}\)}
    {\(1 - 10^{-x}\)}
    {\(10^{-x} - 1\)}
\end{problem}

\begin{problem}
某地 2004 年第一季度应聘和招聘人数排行榜前 5 个行业的情况列表如下:

\begin{center}
\begin{tabular}{|c|c|c|c|c|c|}
\hline
行业名称 & 计算机 & 机械 & 营销 & 物流 & 贸易 \\
\hline
应聘人数 & 215830 & 200250 & 154676 & 74570 & 65280 \\
\hline
\end{tabular}

\medskip

\begin{tabular}{|c|c|c|c|c|c|}
\hline
行业名称 & 计算机 & 营销 & 机械 & 建筑 & 化工 \\
\hline
招聘人数 & 124620 & 102935 & 89115 & 76516 & 70436 \\
\hline
\end{tabular}
\end{center}

若用同一行业中应聘人数与招聘人数比值的大小来衡量该行业的就业情况，则根据表中数据，就业形势一定是（\quad）

\choices
    {计算机行业好于化工行业}
    {建筑行业好于物流行业}
    {机械行业最紧张}
    {营销行业比贸易行业紧张}
\end{problem}

\section{解答题}
\begin{problem}
已知复数 \(z_{1}\) 满足 \((1 + \i)z_{1} = -1 + 5\i, z_{2} = a - 2 - \i\)，其中 \(\i\) 为虚数单位，\(a \in \R\)，若 \(|z_{1} - z_{2}| < |z_{1}|\)，求 \(a\) 的取值范围.
\end{problem}

\begin{problem}
某单位用木料制作如图所示的框架，框架的下部是边长分别为 \(x\)、\(y\)（单位: m）的矩形. 上部是等腰直角三角形. 要求框架围成的总面积 \(8\mathrm{cm}^2\). 问 \(x\)、\(y\) 分别为多少（精确到 \(0.001\mathrm{m}\)）时用料最省?

\bitmapfigure[page=1,pagebox=mediabox,viewport=992 630 1172 738,clip,max width=0.26\linewidth,max totalheight=4.6cm]{source/普通高考/2004/2004上海文.pdf}
\end{problem}

\begin{problem}
记函数 \(f(x) = \sqrt{2 - \frac{x + 3}{x + 1}}\) 的定义域为 \(A\)，\(g(x) = \lg[(x - a - 1)(2a - x)]\) \((a < 1)\) 的定义域为 \(D\).

\begin{enumerate}[itemsep=0pt,topsep=0pt]
\item 求 \(A\);
\item 若 \(B \subseteq A\)，求实数 \(a\) 的取值范围.
\end{enumerate}
\end{problem}

\begin{problem}
如图，直线 \(y = \frac{1}{2}x\) 与抛物线 \(y = \frac{1}{8}x^2 - 4\) 交于 \(A\)、\(B\) 两点，线段 \(AB\) 的垂直平分线与直线 \(y = -5\) 交于 \(Q\) 点.

\bitmapfigure[page=2,pagebox=mediabox,viewport=48 442 408 650,clip,max width=0.36\linewidth,max totalheight=3.3cm]{source/普通高考/2004/2004上海文.pdf}

\begin{enumerate}[itemsep=0pt,topsep=0pt]
\item 求点 \(Q\) 的坐标;
\item 当 \(P\) 为抛物线上位于线段 \(AB\) 下方（含 \(A\)、\(B\)）的动点时，求 \(\triangle OPQ\) 面积的最大值.
\end{enumerate}
\end{problem}

\begin{problem}
如图，\(P-ABC\) 是底面边长为 \(1\) 的正三棱锥，\(D\)、\(E\)、\(F\) 分别为棱长 \(PA\)、\(PB\)、\(PC\) 上的点，截面 \(DEF \parallel\) 底面 \(ABC\)，且棱台 \(DEF-ABC\) 与棱锥 \(P-ABC\) 的棱长和相等.（棱长和是指多面体中所有棱的长度之和）

\bitmapfigure[page=2,pagebox=mediabox,viewport=536 418 760 618,clip,max width=0.28\linewidth,max totalheight=3.2cm]{source/普通高考/2004/2004上海文.pdf}

\begin{enumerate}[itemsep=0pt,topsep=0pt]
\item 证明: \(P-ABC\) 为正四面体;
\item 若 \(PD = \frac{1}{2}PA\)，求二面角 \(D-BC-A\) 的大小.（结果用反三角函数值表示）
\item 设棱台 \(DEF-ABC\) 的体积为 \(V\)，是否存在体积为 \(V\) 且各棱长均相等的直平行六面体，使得它与棱台 \(DEF-ABC\) 有相同的棱长和? 若存在，请具体构造出这样的一个直平行六面体，并给出证明; 若不存在，请说明理由.
\end{enumerate}
\end{problem}

\begin{problem}
设 \(P_{1}(x_{1},y_{1}), P_{2}(x_{2},y_{2}), \cdots, P_{n}(x_{n},y_{n})\) \((n \geqslant 3, n \in \mathbf{N})\) 是二次曲线 \(C\) 上的点，且 \(a_{1} = |OP_{1}|^{2}, a_{2} = |OP_{2}|^{2}, \cdots, a_{n} = |OP_{n}|^{2}\) 构成了一个公差为 \(d\) \((d \neq 0)\) 的等差数列，其中 \(O\) 是坐标原点. 记 \(S_{n} = a_{1} + a_{2} + \cdots + a_{n}\).

\begin{enumerate}[itemsep=0pt,topsep=0pt]
\item 若 \(C\) 的方程为 \(\frac{x^2}{9} - y^2 = 1\)，\(n = 3\). 点 \(P_{1}(3,0)\) 及 \(S_{3} = 162\)，求点 \(P_{3}\) 的坐标;（只需写出一个）
\item 若 \(C\) 的方程为 \(y^{2} = 2px\) \((p \neq 0)\). 点 \(P_{1}(0,0)\)，对于给定的自然数 \(n\)，证明: \((x_{1} + p)^{2}, (x_{2} + p)^{2}, \cdots, (x_{n} + p)^{2}\) 成等差数列;
\item 若 \(C\) 的方程为 \(\frac{x^2}{a^2} + \frac{y^2}{b^2} = 1\) \((a > b > 0)\). 点 \(P_{1}(a,0)\)，对于给定的自然数 \(n\)，当公差 \(d\) 变化时，求 \(S_{n}\) 的最小值.
\end{enumerate}
\end{problem}
